\documentclass[a4paper]{article}
\usepackage{amsmath}
\usepackage{amsfonts}
\usepackage[affil-it]{authblk}
\usepackage[backend=bibtex,style=numeric,sorting=none]{biblatex}
\usepackage{enumitem}
\usepackage{graphicx}
\usepackage{float}

\usepackage{geometry}
\geometry{margin=1.5cm, vmargin={0pt,1cm}}
\setlength{\topmargin}{-1cm}
\setlength{\paperheight}{29.7cm}
\setlength{\textheight}{25.3cm}
\setlength{\parindent}{0pt}

\addbibresource{references.bib}

\begin{document}
% =================================================
\title{Numerical Analysis theoretical homework \# 4}

\author{WangHao 3220103613
  \thanks{Electronic address: \texttt{3220103613@zju.edu.cn}}}
\affil{(Information and Computing Sciences, Class 2201), Zhejiang University }


\date{Due time: \today}

\maketitle
% ============================================
\section*{Exercise 4.33}
$a=1.234*10^4$.
\subsection*{Answer of 4.33(1)}
$b=8.769*10^4$. 
As $a<b$, I swap them. Now we have $a=8.769*10^4$ and $b=1.234*10^4$.

\begin{enumerate}[label=(\roman*)] 
    \item $e_a-e_b=4-4=0$, so $e_c \gets 4$ and $M_b \gets 1.234$.
    \item Perform the addition. $M_c \gets 10.003$.
    \item Normalization. $M_c \gets 1.0003$, $e_c \gets 5$.
    \item Check range. Do nothing.
    \item Round. $M_c \gets 1.000$.
    \item $c \gets 1.000 * 10^5$.
\end{enumerate}

\subsection*{Answer of 4.33(2)}
We have $a=1.234*10^4$ and $b=-5.678*10^0$.

\begin{enumerate}[label=(\roman*)] 
    \item $e_a-e_b=4-0=4$, so $e_c \gets 4$ and $M_b \gets -0.0005678$
    \item Perform the addition. $M_c \gets 1.2334322$.
    \item Normalization. Do nothing.
    \item Check range. Do nothing.
    \item Round. $M_c \gets 1.233$.
    \item $c \gets 1.233 * 10^4$.
\end{enumerate}

\subsection*{Answer of 4.33(3)}
We have $a=1.234*10^4$ and $b=-5.678*10^3$.

\begin{enumerate}[label=(\roman*)] 
    \item $e_a-e_b=4-3=1$, so $e_c \gets 4$ and $M_b \gets -0.5678$
    \item Perform the addition. $M_c \gets 0.6662$.
    \item Normalization. $M_c \gets 6.662$, $e_c \gets 3$.
    \item Check range. Do nothing.
    \item Round. Do nothing.
    \item $c \gets 6.662 * 10^3$.
\end{enumerate}

\section*{Exercise 4.42}

\subsection*{Answer of 4.42}
Consider the FPN system $\mathcal{F}:(10,1,-5,5)$ with $a=1*10^0,b=c=5*10^{-1}$.
Consider the calculation $a+b+c$. precise calculation is $a+b+c=2$.

Consider calculations below: $\text{fl}(\text{fl}(a+b)+c)$, which is not sorted, and $\text{fl}(\text{fl}(b+c)+a)$, which is sorted.

$\text{fl}(\text{fl}(a+b)+c)=1*10^0+5*10^{-1}+5*10^{-1}=2*10^0+5*10^{-1}=3*10^0$. So $\delta>0$.

While $\text{fl}(\text{fl}(b+c)+a)=5*10^{-1}+5*10^{-1}+1*10^0=1*10^0+1*10^0=2*10^0$. So $\delta=0$.

\section*{Exercise 4.43}

\subsection*{Answer of 4.43(1)}
\begin{align*}
    \text{Original formula}&=\text{fl}(\sum_{i=1}^{3} \text{fl}(a_ib_i))\\
    &=(1+\delta_0)\sum_{i=1}^{3} \text{fl}(a_ib_i),\quad \text{where } \delta_0<2\epsilon_u\\
    &=(1+\delta_0)\sum_{i=1}^{3} (a_ib_i(1+\delta_i)),\quad \text{where } \delta_i<\epsilon_u, \forall i \in [3].
\end{align*}
Ignore items of $O(\delta^2)$, we have $$\text{fl}(\sum_{i=1}^{3} (a_ib_i))=(1+\delta)(\sum_{i=1}^{3} \text{fl}(a_ib_i)),\quad \text{where } \delta<3\epsilon_u.$$

\subsection*{Answer of 4.43(2)}
\begin{align*}
    \text{Original formula}&=\text{fl}(\sum_{i=1}^{n} \text{fl}(\prod_{j=1}^{m}a_{i,j}))\\
    &=(1+\delta_0) \sum_{i=1}^{n} \text{fl}(\prod_{j=1}^{m}a_{i,j}),\quad \text{where } \delta_0<(n-1)\epsilon_u\\
    &=(1+\delta_0)\sum_{i=1}^{n} ((\prod_{j=1}^{m}a_{i,j})(\prod_{j=1}^{m-1}(1+\delta_{i,j}))),\quad \text{where } \delta_{i,j}<\epsilon_u, \forall i \in [n], \forall j \in [m-1].
\end{align*}
Use Theorem 4.44 \cite{Trm4_44} and ignore items of $O(\delta^2)$, we have 
$$\text{fl}(\sum_{i=1}^{n} \text{fl}(\prod_{j=1}^{m}a_{i,j}))=(1+\delta)(\sum_{i=1}^{n} (\prod_{j=1}^{m}a_{i,j})),\quad \text{where } \delta<((m-1)+(n-1))\epsilon_u.$$

\section*{Exercise 4.80}

\subsection*{Answer of 4.80}
In Exp 4.79, we have $$\text{cond}_f(x)=\frac{x}{\sin x}.$$
Now,
\begin{align*}
    f_A(x)&=\text{fl}(\frac{\text{fl}(\sin x)}{\text{fl}(1+\text{fl}(\cos x))})\\
    &=(1+\delta_1)\frac{(1+\delta_2)\sin x}{(1+\delta_3)(1+(1+\delta_4)\cos x)},\quad \text{where }\delta_i<\epsilon_u,\forall i \in [4].
\end{align*}
Ignore items of $O(\delta^2)$, we have
$$f_A(x)=\frac{(\sin x)}{(1+(\cos x))}(1+\delta_1+\delta_2-\delta_3-\delta_4\frac{\cos x}{1+\cos x}).$$
Hence we have $\varphi(x)=3+\frac{\cos x}{1+\cos x}$ and 
\begin{align*}
    \text{cond}_A(x) &\leq \frac{\sin x}{x}(3+\frac{\cos x}{1+\cos x})\\
    &\leq 1*(3+1)=4.
\end{align*}
Now we can draw a conclusion that algorithm $f_A$ is well-conditioned.

\section*{I.}

\subsection*{Answer of I}
By doing euclidean division recursively, and connect remainders in reverse order, we have
$$(477)_{10}=(111011101)_2=(1.11011101)_2*2^8.$$

\section*{II.}

\subsection*{Answer of II}
By doing multiplication recursively, and connect in order, we have
$$(\frac{3}{5})_{10}=(0.\overline{1001})_2=(1.001\overline{1001})_2*2^{-1}.$$

\section*{III.}

\subsection*{Proof of III}
We have $e_{x_R}=e$, and $x_R=(1+\epsilon_M)\beta^e$.

However, $e_{x_L} \neq e$, otherwise $m_{x_L}<1$. 
So $e_{x_L} = e-1$, and $x_L=(\beta-\epsilon_M)\beta^{e-1}$.

Hence $x_R-x=\epsilon_M\beta^e$, while $x-x_L=\beta^e-\beta^e+\epsilon_M\beta^{e-1}=\epsilon_M\beta^{e-1}$.

Then we have $x_R-x=\beta(x-x_L)$

\section*{IV.}

\subsection*{Answer of IV(1): find out 2 FPNs}
Denote 2 FPNs as $x_L$ and $x_R$, and $x_L<x<x_R$.

In II, we have $$x=(\frac{3}{5})_{10}=(0.\overline{1001})_2=(1.001\overline{1001})_2*2^{-1}.$$
Rewrite it in normalized FPN way, the first 23 decimal places is $00110011001100110011001$, $e=01111110$ and sign bit is 0.

Thus, $$x_L=00111111000110011001100110011001$$ and $$x_R=00111111000110011001100110011010.$$

\subsection*{Answer of IV(2): fl(x)}
We know $x_R-x_L=\epsilon_M*2^{-1}$. 
And observe that $(x-x_L)*\frac{1}{\epsilon_M}=x*2^{-1}$, it equals to $(x-x_L)=\epsilon_M*x*2^{-1}=\frac{3}{10}\epsilon_M$
Then $x_R-x=(x_R-x_L)-(x-x_L)=\frac{2}{10}\epsilon_M.$

So fl(x)=$x_R$.

\subsection*{Answer of IV(3): roundoff error}
$$R_{rel}=\frac{|\text{fl}(x)-x|}{|x|}=\frac{1}{3}\epsilon_M=\frac{1}{3}*2^{-23}.$$

\section*{V.}

\subsection*{Answer of V}
Denote $$x=(d_0.d_1\ldots d_{p-1})_\beta*\beta^e,\quad \text{thus } m_x=(d_0.d_1\ldots d_{p-1})_\beta$$
and 
$$x'=\arg\max_y |y-x|.$$

$\implies m_{x'}=(d_0.d_1\ldots d_{p-1} \overline{(\beta-1)})_\beta$

Thus, $$\epsilon_u=\epsilon_M*(0.\overline{(\beta-1)})_\beta=\epsilon_M=2^{-23}.$$

\section*{VI.}

\subsection*{Answer of VI}
Apply Taylor expansion to $\cos x$ at $x=0$, we have
$$\cos x=1-\frac{1}{2!}x^2+\frac{1}{4!}x^4-\ldots.$$
Thus, 
$$\cos \frac{1}{4}=1-\frac{1}{2!}(\frac{1}{4})^2+\frac{1}{4!}(\frac{1}{4})^4-\ldots.$$
Because $$1-\frac{1}{2!}(\frac{1}{4})^2=1.000000-0.000010=0.111100,$$
and we can observe that the remainders will not influence the first 6 decimal places.

Then $$1-\cos x=1-0.111100\ldots=0.00001\ldots,$$ 5 bits of precision lost.

\section*{VII.}

\subsection*{Answer of VII: way 1}
Using Taylor expansion at $x=0$ before calculation, we have
\begin{align*}
    1-\cos x&=1-(1-\frac{1}{2!}x^2+\frac{1}{4!}x^4-\ldots)\\
    &=\frac{1}{2!}x^2-\frac{1}{4!}x^4+\ldots
\end{align*}
And calculate it. Now $\delta<\epsilon_u$.

\subsection*{Answer of VII: way 2}
Using trigonometric identity before calculation, we have
\begin{align*}
    1-\cos x&=1-(1-2\sin^2 \frac{x}{2})\\
    &=2\sin^2 \frac{x}{2}
\end{align*}
And calculate it. Now $\delta<\epsilon_u$.

\section*{VIII.}
$$\text{cond}_f(x)=|\frac{xf'(x)}{f(x)}|.$$

\subsection*{Answer of VIII(1)}
Substitute $f(x)$ into formula, we have
$$\text{cond}_f(x)=|\frac{\alpha x}{x-1}|.$$
It is large when $x \to 1$.

\subsection*{Answer of VIII(2)}
Substitute $f(x)$ into formula, we have
$$\text{cond}_f(x)=\frac{1}{|\ln x|}.$$
It is large when $x \to 1$.

\subsection*{Answer of VIII(3)}
Substitute $f(x)$ into formula, we have
$$\text{cond}_f(x)=|x|.$$
It is large when $x \to \infty$.

\subsection*{Answer of VIII(4)}
Substitute $f(x)$ into formula, we have
$$\text{cond}_f(x)=|\frac{x}{\sqrt{1-x^2}\arccos x}|.$$
It is large when $x \to -1^+,1^-$.

\section*{IX.}

\subsection*{Proof of IX(1)}
$$\text{cond}_f(x)=|\frac{xe^{-x}}{1-e^{-x}}|=\frac{x}{e^x-1}\leq\frac{x}{x}=1.$$

\subsection*{Answer of IX(2)}
\begin{align*}
    f_A(x)&=\text{fl}(1-\text{fl}(e^{-x}))\\
    &=(1+\delta_2)(1-e^{-x}(1+\delta_1)),\quad \text{where }\delta_1,\delta_2<\epsilon_u\\
    &=(1-e^{-x})(1+\delta_2)(\frac{1-e^{-x}(1+\delta_1)}{1-e^{-x}})\\
    &=(1-e^{-x})(1+\delta_2)(1-\frac{e^{-x}}{1-e^{-x}}\delta_1)\\
    &=(1-e^{-x})(1+\delta_2-\frac{e^{-x}}{1-e^{-x}}\delta_1-\frac{e^{-x}}{1-e^{-x}}\delta_1\delta_2).
\end{align*}

$\implies \varphi(x)=1+\frac{e^{-x}}{1-e^{-x}}(1+\epsilon_u)$

$\implies \text{cond}_A(x)<\frac{1+\frac{e^{-x}}{1-e^{-x}}(1+\epsilon_u)}{\frac{xe^{-x}}{1-e^{-x}}}=\frac{e^x+\epsilon_u}{x}.$

\subsection*{Answer of IX(3)}

\begin{figure}[H]
    \begin{minipage}{0.4\textwidth}
        \centering
        \includegraphics[width=\textwidth]{IX(3.1).pdf}
        \caption{Image of $\text{cond}_f(x)$}
    \end{minipage}
    \hfill
    \begin{minipage}{0.4\textwidth}
        \centering
        \includegraphics[width=\textwidth]{IX(3.2).pdf}
        \caption{Image of upper bound of $\text{cond}_A(x)$}
    \end{minipage}
\end{figure}

We can find that $\text{cond}_A(x) \to \infty\quad(x \to 0)$. It have no upper bound.
This is because calculation of $\text{fl}(\text{fl}(x)-\text{fl}(y))$ is not accurate when $x \approx y$.

While when $x$ is close to 1, $\text{cond}_A(x)$ is acceptable.\\
This is because the calculation generates acceptable error and $\text{cond}_f(x)$ is also acceptable when $x$ is close to 1.

\section*{X.}

\subsection*{Proof of X(1)}
By property of singular value, $\|A\|_2=\sigma_{max}(A)$ and $\|A^{-1}\|_2=\sigma_{min}(A)^{-1}$.\\
Then $$\text{cond}_2A=\|A\|_2\|A^{-1}\|_2=\frac{\sigma_{max}(A)}{\sigma_{min}(A)}$$

\subsection*{Proof of X(2)}
As A is normal, we have
$$A=U^T \Sigma U, \quad \text{where } \Sigma=\mathrm{diag}(\lambda_{\max}(A), \ldots, \lambda_{\min}(A)).$$
We also have
$$A^TA=U^T \Sigma^2 U, \quad \text{where } \Sigma=\mathrm{diag}(\lambda_{\max}^2(A), \ldots, \lambda_{\min}^2(A)).$$
This is also a SVD of $A^TA$.\\
Then,
\begin{align*}
    \text{cond}_2A&=\frac{\sigma_{max}(A)}{\sigma_{min}(A)}\\
    &=\sqrt{\frac{\sigma_{max}^2(A)}{\sigma_{min}^2(A)}}\\
    &=\sqrt{\frac{\sigma_{max}(A^TA)}{\sigma_{min}(A^TA)}}\\
    &=\sqrt{\frac{\lambda_{max}(A^TA)}{\lambda_{min}(A^TA)}}\\
    &=\sqrt{\frac{|\lambda_{max}^2(A)|}{|\lambda_{min}^2(A)|}}\\
    &=\frac{|\lambda_{max}(A)|}{|\lambda_{min}(A)|}
\end{align*}

\subsection*{Proof of X(3)}
A is unitary\\
$\implies \lambda_{max}(A)=\lambda_{min}(A)=1$\\
$\implies \text{cond}_2A=1$

\section*{XI.}

\subsection*{Answer of XI(1)}
Denote $\vec{a}=(a_0,a_1,\ldots,a_{n-1})$.
$$\text{cond}_f(\vec{a})=\frac{\|\vec{a}\|\|\nabla f\|}{|r|},\quad \text{where } (\nabla f)_{i+1}=\frac{dr}{da_i}.$$
By Lemma 4.63 \cite{Lemma4_63}, $\frac{dr}{da_i}=-\frac{r^i}{q'(r)}$.
\begin{align*}
    \implies \text{cond}_f(\vec{a})&=\frac{(\sum_{i=0}^{n-1}|a_i|)(\sum_{i=0}^{n-1}|\frac{r^i}{q'(r)}|)}{r}\\
    &=\frac{(\sum_{i=0}^{n-1}|a_i|)(\sum_{i=0}^{n-1}|r^i|)}{q'(r)}.
\end{align*}

\subsection*{Answer of XI(2)}
Apply on Wilkinson example, we have
$$\text{cond}_f(\vec{a})=\frac{p+1}{p-1}(p^p-1).$$
With a disturbance of $\epsilon x^p$, we have
$$\Delta r=r*\text{cond}_A(\vec{a})*\frac{\|\Delta \vec{a}\|}{\|\vec{a}\|}=\frac{p^p-1}{(p-1)!(p-1)},$$
which is different of result of Exp 4.64 \cite{Exp4_64} by only a constant.

This implies that the condition number can evaluate the stability of the problem.

\section*{XII.}

\subsection*{Answer of XII}
Below is the contradiction.

Assume $p=2$. Precision of register is 4 bits.

Consider division of $$x=\frac{1}{(1.1)_2}.$$
In register, the result is $$.1010$$ After round to even, we have $$x=1.0*2^{-1}$$
Now $$\delta=|\frac{\frac{1}{(1.1)_2}-1.0*2^{-1}}{\frac{1}{(1.1)_2}}|=(0.01)_2,$$
and $$\epsilon_u=0.5*2^{-1}=(0.01)_2=\delta.$$

\section*{XIII.}

\subsection*{Answer of XIII}
We cannot guarantee that. Below is the explaination.

As$$x \in [128,129] \subset [2^7,2^8],$$
then we have $e_x=7$.

So $$\max_x E_{abs}(x)=\frac{\epsilon_u}{1+\epsilon_u}*2^7\approx2^{-17}>10^{-6}.$$
Thus we cannot guarantee the absolute error of x $<10^{-6}$.

% ============================================
\section*{XIV.}

\subsection*{Answer of XIV}
Denote 2 adjacent points as $\overline{x}$ and $\overline{x}+\Delta x$.

We can write 4 equations about coefficients of cubic splines in format of matrix below.

Denote
$$
M=\begin{bmatrix}
    1 & 0 & 0 & 0 \\
    0 & 1 & 0 & 0 \\
    1 & \Delta x & \Delta x^2 & \Delta x^3 \\
    0 & 1 & 2\Delta x & 3\Delta x^2
    \end{bmatrix}
    .
$$

Then we can solve $a_0,a_1,a_2,a_3$ from

$$
M
\begin{bmatrix}
a_0 \\ a_1 \\ a_2 \\ a_3
\end{bmatrix}
=
\begin{bmatrix}
f(\overline{x}) \\ f'(\overline{x}) \\ f(\overline{x} + \Delta x) \\ f'(\overline{x} + \Delta x)
\end{bmatrix}
.
$$

We can observe that 
$$\lim_{\Delta x \to 0}M \to 
\begin{bmatrix}
    1 & 0 & 0 & 0 \\
    0 & 1 & 0 & 0 \\
    1 & 0 & 0 & 0 \\
    0 & 1 & 0 & 0
    \end{bmatrix}
    ,
$$
which is a singular matrix.

Thus $$\sigma_{min}(A) \to 0 \quad (\Delta x \to 0),$$
then $$\text{cond}_2A=\frac{\sigma_{max}(A)}{\sigma_{min}(A)} \to \infty \quad (\Delta x \to 0).$$ 

Which means the error of $a_0,a_1,a_2,a_3$ can be very large. Then the curve is inaccurate.

\printbibliography

\end{document}